\usepackage[scr=rsfs,cal=euler]{mathalfa}
\renewcommand{\thetable}{\arabic{table}}
\multlinegap0pt

\hyphenation{every-thing recall table}

\usepackage[capitalise]{cleveref}
\crefdefaultlabelformat{#2\textup{#1}#3}

\makeatletter
\def\cref@thmoptarg[#1]#2#3#4{%
    \ifhmode\unskip\unskip\par\fi%
    \normalfont%
    \trivlist%
    \let\thmheadnl\relax%
    \let\thm@swap\@gobble%
    \thm@notefont{\fontseries\mddefault\upshape}%
    \thm@headpunct{.}% add period after heading
    \thm@headsep 5\p@ plus\p@ minus\p@\relax%
    \thm@space@setup%
    #2% style overrides
    \@topsep \thm@preskip               % used by thm head
    \@topsepadd \thm@postskip           % used by \@endparenv
    \def\@tempa{#3}\ifx\@empty\@tempa%
      \def\@tempa{\@oparg{\@begintheorem{#4}{}}[]}%
    \else%
      \refstepcounter[#1]{#3}%  <<< cleveref modification
      \@namedef{cref@#3@alias}{#1}% added
      \def\@tempa{\@oparg{\@begintheorem{#4}{\csname the#3\endcsname}}[]}%
    \fi%
    \@tempa}%
\makeatother

\let\oldbigsqcup\bigsqcup
\renewcommand{\bigsqcup}{\mathop{\textstyle\oldbigsqcup}\displaylimits}
\newcommand\mto{\mathchoice{\longmapsto}{\mapsto}{\mapsto}{\mapsto}}
\newcommand{\sbullet}{{\scriptscriptstyle\bullet}}
\newcommand{\csbullet}{{\raisebox{1pt}{$\sbullet$}}}

\let\epsilon\varepsilon
\let\ra\rightarrow

\usepackage{booktabs}
\usepackage{colonequals}
\usepackage{tikz-cd}

\theoremstyle{plain}
\newtheorem{theorem}{Theorem}[section]
\newtheorem{corollary}[theorem]{Corollary}
\newtheorem{lemma}[theorem]{Lemma}
\newtheorem{proposition}[theorem]{Proposition}

\newtheorem{alphatheorem}{Theorem}
\newtheorem{alphaconjecture}[alphatheorem]{Conjecture}
\newtheorem{alphacorollary}[alphatheorem]{Corollary}
\newtheorem{alphaproposition}[alphatheorem]{Proposition}

\theoremstyle{definition}
\newtheorem{definition}[theorem]{Definition}
\newtheorem{remark}[theorem]{Remark}
\newtheorem{example}[theorem]{Example}

\renewcommand{\thealphatheorem}{\Alph{alphatheorem}}
\renewcommand{\thealphaconjecture}{\Alph{alphatheorem}}
\renewcommand{\thealphacorollary}{\Alph{alphatheorem}}
\renewcommand{\thealphaproposition}{\Alph{alphaproposition}}

\crefname{alphatheorem}{Theorem}{Theorems}
\crefname{alphaconjecture}{Conjecture}{Conjectures}
\crefname{alphacorollary}{Corollary}{Corollaries}
\crefname{alphaproposition}{Proposition}{Propositions}

\DeclareMathOperator\Br{Br}
\DeclareMathOperator\characters{X}
\DeclareMathOperator\derived{\mathbf{D}}
\DeclareMathOperator\dimvect{\underline{\operatorname{dim}}}
\DeclareMathOperator\End{End}
\DeclareMathOperator\Ext{Ext}
\DeclareMathOperator\GL{GL}
\DeclareMathOperator\HH{H}
\DeclareMathOperator\hochschild{HH}
\DeclareMathOperator\Hilb{Hilb}
\DeclareMathOperator\Hom{Hom}
\DeclareMathOperator\sheafHom{\mathcal{H}om}
\DeclareMathOperator\RsheafHom{\mathbf{R}\mathcal{H}om}
\DeclareMathOperator\Mat{Mat}
\DeclareMathOperator\PG{PG}
\DeclareMathOperator\Pic{Pic}
\DeclareMathOperator\Rep{R}
\DeclareMathOperator\rk{rk}
\DeclareMathOperator\Spec{Spec}
\DeclareMathOperator\source{s}
\DeclareMathOperator\target{t}
\DeclareMathOperator\weight{wt}

\mathchardef\mhyphen="2D

\newcommand\blade[1][]{\Sigma_{#1}}
\newcommand\bounded{\ensuremath{\mathrm{b}}}
\newcommand\can{\ensuremath{\mathrm{can}}}
\newcommand\cocharacters{\ensuremath{\mathrm{X}_*}}
\newcommand\dash{\nobreakdash-\hspace{0pt}}
\newcommand\diff{\mathrm{d} \kern 0pt}
\newcommand\field{\mathbf{k}}
\newcommand\Gm{\ensuremath{\mathbf{G}_{\mathrm{m}}}}
\newcommand\group[1][\tuple{d}]{\mathrm{G}_{#1}}
\newcommand\HN{\ensuremath{{\scriptscriptstyle\mathrm{HN}}}}
\newcommand\limitset[1][\lambda]{Z_{#1}}
\newcommand\moduli{\ensuremath{\mathrm{M}}}
\newcommand\normal{\ensuremath{\mathrm{N}}}
\newcommand\semisimple{\ensuremath{\mathrm{ssimp}}}
\newcommand\semistable{\ensuremath{\mhyphen\mathrm{sst}}}
\newcommand\stable{\ensuremath{\mhyphen\mathrm{st}}}
\newcommand\tangent{\ensuremath{\mathrm{T}}}
\newcommand\tuple[1]{\ensuremath{{\boldsymbol{#1}}}}

\DeclareDocumentCommand\modulistack{om}{\IfNoValueTF{#1}{\mathcal{M}{(#2)}}{\mathcal{M}^{#1}(#2)}}
\DeclareDocumentCommand\modulispace{om}{\IfNoValueTF{#1}{\mathrm{M}{(#2)}}{\mathrm{M}^{#1}(#2)}}
\DeclareDocumentCommand\repspace{om}{\IfNoValueTF{#1}{\mathrm{R}{(#2)}}{\mathrm{R}^{#1}(#2)}}

\newif\ifcode

